\section{Architecture du système}

\subsection{Vue d'ensemble}

L'architecture de \AUSSI{} repose sur la séparation entre le moteur de
recherche générique et les problèmes particuliers. Le moteur reçoit un objet
respectant l'interface \code{Problem}, une stratégie de recherche et un mode
de gestion des états.

La fonction centrale du système est définie comme suit :

\begin{center}
\code{SEARCH(problem, strategy, mode)}
\end{center}

\begin{figure}[H]
\centering
\begin{tikzpicture}[node distance=1.1cm and 1.3cm]

\node[module] (ui) {Interface utilisateur\\CLI};
\node[module, right=of ui] (search) {Moteur générique\\\code{SEARCH()}};
\node[module, right=of search] (metrics) {Résultats et\\métriques};

\node[component, below=of search] (strategy) {Stratégie\\BFS / UCS / Greedy / A*};
\node[component, below=of ui] (mode) {Mode\\Tree / Graph};
\node[component, below=of metrics] (problem) {Problèmes\\Labyrinthe, navigation,\\N-Puzzle, arithmétique};

\draw[arrow] (ui) -- (search);
\draw[arrow] (search) -- (metrics);
\draw[arrow] (strategy) -- (search);
\draw[arrow] (mode) -- (search);
\draw[arrow] (problem) -- (search);

\end{tikzpicture}
\caption{Architecture générale de \AUSSI{}.}
\label{fig:architecture}
\end{figure}

\subsection{Interface \code{Problem}}

Tous les domaines respectent une même interface. Cette interface encapsule les
parties spécifiques à un problème de recherche.

\begin{lstlisting}[caption={Interface conceptuelle d'un problème.}, label={lst:problem-interface}]
class Problem:
    def initial_state(self):
        pass

    def is_goal(self, state) -> bool:
        pass

    def actions(self, state):
        pass

    def result(self, state, action):
        pass

    def step_cost(self, state, action, next_state) -> float:
        pass

    def heuristic(self, state) -> float:
        return 0.0

    def state_key(self, state):
        return state
\end{lstlisting}

La méthode \code{state\_key} fournit une représentation immuable et comparable
d'un état. Elle est utilisée par le mode \GraphSearch{} pour détecter les
états déjà rencontrés et conserver le meilleur coût connu.

\subsection{Noeud de recherche}

Un noeud de recherche contient au minimum l'état courant, un pointeur vers son
parent, l'action ayant produit l'état, le coût cumulé $g(n)$ et la profondeur.
Cette structure permet de reconstruire la solution après l'atteinte d'un état
but.

\begin{lstlisting}[caption={Structure conceptuelle d'un noeud de recherche.}]
@dataclass
class Node:
    state: Any
    parent: Node | None
    action: Any | None
    path_cost: float
    depth: int

    def solution(self):
        actions = []
        node = self
        while node.parent is not None:
            actions.append(node.action)
            node = node.parent
        return list(reversed(actions))
\end{lstlisting}

\subsection{Résultat standardisé}

Le moteur retourne un objet de résultat commun à tous les problèmes. Celui-ci
contient notamment le statut de résolution, le chemin, le coût, la profondeur,
le nombre de noeuds générés, le nombre de noeuds expansés, la taille maximale
de la frontière et le temps d'exécution.